\documentclass[10pt]{article}

\usepackage[utf8]{inputenc}

\usepackage[T1]{fontenc}

\usepackage{amsmath}

\usepackage{amsfonts}

\usepackage{amssymb}

\usepackage[version=4]{mhchem}

\usepackage{stmaryrd}

\usepackage{bbold}



\begin{document}

ков» с плотностями заряда $e \psi_{n}^{*}\left(r_{1}\right) \psi_{m}\left(r_{1}\right)$ и $e \psi_{m}^{*}\left(r_{2}\right) \psi_{n}\left(r_{2}\right)$ (звездочка означает комплексное сопряжение), то есть когда каждый из электронов находится одновременно в состояниях $\psi_{n}$ и $\psi_{m}$ (что бессмысленно с точки зрения классич. физики). Из (3) следует, что полная энергия пара- и ортогелия с электронами в аналогичных состояниях отличается на величину $2 A$. Таким образом, хотя непосредственно спиновое взаимодействие мало и не учитывается, тождественность двух электронов в атоме гелия приводит к тому, что энергия системы оказывается зависящей от полного спина системы, как если бы между частицами существовало дополнительное, обменное взаимодействие. Очевидно, что О. в. в данном случае является частью кулоновского взаимодействия электронов и явным образом выступает при приближенном рассмотрении квантовомеханич. системы, когда волновая функция всей системы выражается через волновые функции отдельных частиц (в частности, в приближении Хартри - Фока).



О. в. эффективно проявляется в тех случаях, когда «перекрываются» волновые функции отдельных частиц системы, то есть когда существуют области пространства, в К-рых с заметной вероятностью может находиться частица в различных состояниях движения. Это видно из выражения для обменного интеграла $A$ : если степень перекрытия состояний $\psi_{n}^{*}(r)$ и $\psi(r)$ незначительна, то величина $A$ очень мала.



Характер О. в. различен для фермионов и для бозонов. Для фермионов О. в. является следствием Паули принципа, препятствующего сближению тождественных частиц с одинаковым направлением спинов, и эффективно проявляется как отталкивание их друг от друга.



В системе тождественных бозонов О. в., напротив, имеет характер взаимного притяжения частиц. Рассмотрение систем из большего числа одинаковых частиц производится на основе Ферми - Дирака статистики для фермионов и Бозе - Эйнштейна статистики для бозонов.



Если взаимодействующие тождественные частицы находятся во внешнем поле, напр. в кулоновском поле ядра, то существование определенной симметрии волновой функции и соответственно определенной корреляции движения частиц влияет на их энергию в этом поле, что также является обменным эффектом. Обычно (в атоме, молекуле, кристалле) это О. в. вносит вклад обратного знака по сравнению с вкладом О. в. частиц друг с другом. Поэтому суммарный обменный эффект может как понижать, так и повышать полную энергию взаимодействия в системе. Энергетич. выгодность или невыгодность состояния с параллельными спинами фермионов, в частности электронов, зависит от относительных величин этих вкладов. Так, в ферромагнетике (аналогично рассмотренному атому гелия) более низкой энергией обладает состояние, в к-ром спины (и магнитные моменты) электронов в незаполненных оболочках соседних атомов параллельны; в этом случае благодаря О. в. возникает спонтанная намагниченность. Напротив, в молекулах с ковалентной химич. связью, напр. в молекуле $\mathrm{H}_{2}$, энергетически выгодно состояние, в к-ром спины валентных электронов соединяющихся атомов антипараллельны.



О. в. объясняет, таким образом, закономерности атомной и молекулярной спектроскопии, химич. связь в молекулах, ферромагнетизм (и антиферромагнетизм), а также другие специфич. явления в системах одинаковых частиц.



Д. А. Киржнии, С. С. Герштейн.



ОБМОТКА TOPA - траектория постоянного уравнения $\dot{\varphi}=\omega$ на $n$-мерном торе $T^{n}$ с угловыми координатами $\varphi \bmod 2 \pi$, где $\varphi=\left(\varphi_{1}, \ldots, \varphi_{n}\right), \omega=\left(\omega_{1}, \ldots, \omega_{n}\right)=$ const $\neq 0$. Иначе говоря, О.т.- это траектория условно периодич. движения на торе, то есть движения, задаваемого постоянным векторным полем. Свойства всех таких траекторий на торе одинаковы и определяются вектором частот $\omega$. Более



396 OБMEHHOE точно - максимальным числом $r$ линейно независимых целочисленных соотношений



\[

(k, \omega)=\sum k_{i} \omega_{i}=0

\]



между частотами $\omega_{i}$, где $k=\left(k_{1}, \ldots, k_{n}\right)$ принадлежит решетке $\mathbb{Z}^{n}$ целочисленных векторов. Соотношение $(k, \omega)=0$ называется резонансом, число $|k|=\sum\left|k_{i}\right|$ - порядком резонанса, a $r-$ кратностью резонанса вектора частот $\omega$.



Замыкание О. т. является $(n-r)$-мерным тором, совпадающим с $T^{n}$ в случае $r=0$. Так что одна О.т. полностью «обматывает» весь тор $T^{n}$ только в нерезонансном случае, но это общий случай. В случае максимальной кратности резонанса $(r=n-1)$ траектории, напротив, периодичны, а вектор $\omega$ коллинеарен целочисленному: $\omega=c K$, где $K \in \mathbb{Z}^{n} \backslash\{0\}$.



Инвариантные торы с условно периодич. движениями возникают в самых различных системах дифференциальных уравнений. Особенно часто они встречаются в гамильтоновых системах, причем, почти исключительно, семействами. Если фазовое пространство сплошь расслоено на такие торы, то систему называют (то р и ч е с к и) и н т е г р и р у е м о й. В фазовом пространстве гамильтоновой системы в такой ситуации могут быть введены канонич. переменные $I, \varphi \bmod 2 \pi$, в к-рых функция Гамильтона $H$ принимает вид $H=H(I)$ и к-рые называются пе ре мен ны е «де й с твие - уголу.



Нерезонансная система $\dot{\varphi}=\omega$ является простейшим примером эргодич. динамич. системы: временное среднее функции в силу системы равно пространственному среднему. Свойства вектора $\omega$ определяют свойства не только постоянных систем, но и их возмущений, напр. возмущений интегрируемых систем. В частности, при «включении» возмущения нек-рые нерезонансные торы невозмущенной гамильтоновой системы не исчезают, а лишь слегка деформируются, образуя в фазовом пространстве большое по мере множество (см. KAM-теория). Движение по этим торам, называемым к о л м о г о р о в с к и м и, будет условно-периодическим, а траектории - нерезонансными О. т.



Грубое объяснение этого факта - вышеупомянутая эргодичность невозмущенной системы на торах $I=$ const с нерезонансными частотами



\[

\omega(I)=\partial H_{0}(I) / \partial I,

\]



где $H=H_{0}(I)+\varepsilon H_{1}(I, \varphi), \varepsilon \ll 1,-$ функция Гамильтона системы. Эргодичность позволяет оценить скорость изменения функций на решениях через пространственное среднее по тору скорости изменения этих функций; а пространственное среднее вектор-функщии $I=-\partial H / \partial \varphi$, отвечающей за движение поперек слоев - торов $I=$ const, по каждому такому тору равно нулю. В действительности здесь, как и при исследовании других аналогичных случаев, важна не просто обусловленная нерезонансностью эргодичность, то есть неравенство нулю выражений $(k, \omega)$, но и удаленность модулей $|(k, \omega)|$ при $k \in \mathbb{Z}^{n} \backslash\{0\}$ от нуля. Выражения $(k, \omega)$ называются малыми знаменателями.



Лит.: [1] Арнольд В. И., Обыкновенные дифференциальные уравнения, 3 изд., М., 1984; [2] е го же, Математические методы классической механики, 3 изд., М., 1989; [3] Нехороше в Н.Н., «Труды Моск. математич. об-ва», 1972, т. 26, с. 181-98.



H. H. Hexopoues.



ОБОБЩЕННАЯ ГИПЕРТЕОМЕТРИЧЕСКАЯ ФУНКЦиЯ - см. Гипергеометрическая функция.



ОБОБЩЕННАЯ ФУНКЦИЯ - математическое понятие, расширяющее классическое понятие функции. На О.ф. распространяются основные линейные операции математич. анализа, причем формулировки приобретают простую замкнутую форму. С физич. точки зрения О. Ф. отражают невозможность абсолютно точного измерения значения физич. величины в точке, поскольку всякий измерительный прибор имеет свою разрешающую способность и показывает нек-рый усредненный «размазанный» результат.





\end{document}